% \documentclass{report}

% \usepackage{graphicx}
% \usepackage{dcolumn}
% \usepackage{bm}
% \usepackage{amsfonts, amssymb, amsmath}
% \usepackage{braket}

%\newcommand{\ket}[1]{\textstyle \left| #1 \right\rangle \displaystyle}
%\newcommand{\bra}[1]{\textstyle \left\langle #1 \right| \displaystyle}
%\newcommand{\braket}[2]{\textstyle \left\langle #1 \left|  #2 \right\rangle\right. \displaystyle}
%\newcommand{\im}{\mathrm{i}}
%\newcommand{\e}{\mathrm{e}}
%\newcommand{\pwrt}[2]{\frac{\partial #1}{\partial #2}}
%\newcommand{\diff}{\mathrm{d}}
%
%\begin{document}

\subsection{Introducing variational anisotropy.}
\label{subsection:variational-anisotropy}

We previously stated in Eq. (\ref{basis-1}) that our variational subbases will
go as 
\begin{equation}
  \Phi_{klm;\eta;\bm{r}}
  \propto
  F_{kl;\eta;r}^{\left(W\right)}
  Y_{lm;\chi\phi}.
\end{equation}
To help capture the anisotropic quality of the wavefunctions introduced by the
anisotropic effective mass, we introduce a scaling along the particular
crystallographic axis. Without loss of generality, we can introduce a scaling
$z\rightarrow \lambda z$ with $\lambda$ playing the role of a second variational
parameter for the subbasis alongside $\eta$. In spherical coordinates, this
scaling takes the form
\begin{align}
  r
  \rightarrow &
  r \vartheta_{\lambda\chi},
  \\
  \chi
  \rightarrow &
  \frac{\lambda\chi}{\vartheta_{\lambda\chi}},
\end{align}
where 
\begin{align}
  \vartheta_{\lambda\chi}
  = &
  \sqrt{ 1 + \left(\lambda^2 - 1\right)\chi^2 },
\end{align}
with $\chi = z/r = \cos\left(\theta\right)$.

This argument can be specialized for the $x$- and $y$-axes as well. To
accomplish this we can define a set of equivalent spherical coordinates
$(\chi_{q},\phi_{q})$ defined in the standard way, but taking another axis as
the zenith. To wit, we have
\begin{align}
  \chi_{x} &\equiv \frac{x}{r} = \sqrt{1-\chi^{2}}\cos\left(\phi\right) \\
  \chi_{y} &\equiv \frac{y}{r} = \sqrt{1-\chi^{2}}\sin\left(\phi\right) \\
  \chi_{z} &\equiv \frac{z}{r} = \chi
\end{align}
\begin{align}
  \phi_{x} &\equiv \atantwo\left(\chi_{z},\chi_{y}\right)\\
  \phi_{y} &\equiv \atantwo\left(\chi_{x},\chi_{z}\right)\\
  \phi_{z} &\equiv \atantwo\left(\chi_{y},\chi_{x}\right) = \phi
\end{align}
where $\atantwo\left(b,a\right) = \text{Arg}\left(a+\im{b}\right)$.

We will need to examine the implications of this introduced anisotropy upon the
symmetry of the wavefunctions, particularly in the context of the symmetry of
the silicon donor. For now, we can state our basis states with both radial and
anisotropic scaling variational parameters as
\begin{equation}
  \Phi_{klm;\eta\lambda;\bm{r}}^{(Wq)}
  =
  F_{lk;\eta\lambda;r\chi_{q}}^{\left(W\right)}
  Y_{lm;\lambda;\chi_{q}\phi_{q}},
  \label{scaled_wf}
\end{equation}
where
\begin{align}
  F_{lk;\eta\lambda;r\chi_{q}}^{\left(W\right)}
  & = 
  F_{lk;\eta_{\lambda;\chi_{q}};r}^{\left(W\right)}
  \quad\quad
  \eta_{\lambda;\chi_{q}}
  =
  \frac{\eta}{\vartheta_{\lambda;\chi_{q}}}
  \label{scaled_wf:radial_part}
  \\
  Y_{lm;\lambda;\chi_{q}\phi_{q}}
  & =
  \sqrt{ \frac{\lambda}{\vartheta_{\lambda;\chi_{q}}^{3}} }
  Y_{
    lm;
    \chi_{\lambda}^{\left(q\right)}
    \phi_{q}
  }
  \quad\quad
  \chi_{\lambda}^{\left(q\right)}
  =
  \frac{\lambda\chi_{q}}{\vartheta_{\lambda;\chi_{q}}}
  \label{scaled_wf:aniso_part}
\end{align}
The additional factors in front of the $F$ and $Y$ functions help to maintain a
degree of independent normalization between the parts of the basis state, and
makes the scaling transformation of some integrals within this project easier,
aiding analytical simplification and helping numerical stability during
computation.